\documentclass[12pt,a4paper]{article}
\usepackage[utf8]{inputenc}
\usepackage[brazil]{babel}
\usepackage[T1]{fontenc}
% Fontes vetoriais; sem isso o T1 cai nas fontes em bitmap.
\usepackage{lmodern}
\usepackage{geometry}
\usepackage{graphicx}
\usepackage{hyperref}
\usepackage{listings}
\usepackage{xcolor}
\usepackage{amsmath}
\usepackage{amssymb}

\geometry{a4paper, left=3cm, right=2cm, top=3cm, bottom=2cm}

\definecolor{codegreen}{rgb}{0,0.6,0}
\definecolor{codegray}{rgb}{0.5,0.5,0.5}
\definecolor{codepurple}{rgb}{0.58,0,0.82}
\definecolor{backcolour}{rgb}{0.95,0.95,0.92}

\lstdefinestyle{mystyle}{
    backgroundcolor=\color{backcolour},
    commentstyle=\color{codegreen},
    keywordstyle=\color{magenta},
    numberstyle=\tiny\color{codegray},
    stringstyle=\color{codepurple},
    basicstyle=\ttfamily\footnotesize,
    breakatwhitespace=false,
    breaklines=true,
    captionpos=b,
    keepspaces=true,
    numbers=left,
    numbersep=5pt,
    showspaces=false,
    showstringspaces=false,
    showtabs=false,
    tabsize=2
}

\lstset{style=mystyle}

% Codigo e saidas curtas, sem numero de linha.
\lstdefinestyle{gram}{style=mystyle, numbers=none, basicstyle=\ttfamily\scriptsize}

% Evita linhas estouradas por trechos longos em \texttt.
\setlength{\emergencystretch}{1em}

% Nomes de tipo nas regras.
\newcommand{\ty}[1]{\mathsf{#1}}

\title{Relatório de Projeto: Compilador MiniPython\\[4pt]
\large Trabalho Prático 2: Verificação de Tipos e Interpretação}
\author{
    Brenda Gabrielle Alves Nascimento \\
    Matrícula: 24.1.4011 \\
    BCC328 - Construção de Compiladores I - DECOM/UFOP
}
\date{\today}

\begin{document}

\maketitle

\begin{abstract}
Segunda etapa do compilador MiniPython, em Haskell: um verificador de
tipos estático, baseado em unificação e com subtipagem entre classes, e
um interpretador que executa a AST dos programas bem tipados. Os dois
entram no executável \texttt{mpyc} pelas opções \texttt{-{}-check} e
\texttt{-{}-run}.
\end{abstract}

\section{Introdução}

O TP1 deixou prontos o analisador léxico e os três parsers, que
produzem a mesma AST. Este trabalho acrescenta as duas fases seguintes:
a verificação de tipos, que confere as regras da seção 10 da
especificação e acumula todos os erros encontrados, e a interpretação,
que executa a AST quando a verificação passa. Com \texttt{-{}-check} e
\texttt{-{}-run}, o programa é lido pelo analisador descendente
recursivo; qualquer um dos três serviria, já que a AST é a mesma. O
código novo está em \texttt{src/MiniPython/Semantic} (tipos e
verificador) e \texttt{src/MiniPython/Interpreter} (valores e
interpretador).

\section{Metodologia}

\subsection{Sistema de tipos para MiniPython}

Os tipos do verificador (\texttt{Semantic/Types.hs}) são separados das
anotações da AST, que carregam posição:

\begin{lstlisting}[style=gram]
data Ty = TInt | TFloat | TBool | TStr | TNone
        | TList Ty | TFun [Ty] Ty | TClass Name
        | TGen Name     -- variavel de tipo escrita na anotacao (a, b, ...)
        | TMeta Int     -- incognita criada pela inferencia (?0, ?1, ...)
\end{lstlisting}

O ambiente de tipos é uma pilha de escopos, cada um um mapa de nomes
para \texttt{BVar} (variável ou parâmetro, com o seu tipo),
\texttt{BFun} (função declarada com \texttt{def}, com um esquema de
tipo) ou \texttt{BClass}. Cada bloco empilha um escopo, descartado na
saída, e o escopo global fica à parte. O corpo de uma função começa com
um escopo só com os parâmetros e não enxerga os escopos locais de quem
a definiu, porque \texttt{def} não forma closure (regra 7 da seção 10);
o global continua visível. O ambiente de classes guarda, para cada
classe, o pai, os campos e os métodos, já somados aos herdados, cada
método com a classe que o definiu.

As regras principais estão abaixo, com $\Gamma$ o ambiente,
$\sigma \le \tau$ a subtipagem e $\Leftarrow$ a verificação contra um
tipo esperado, usada na lambda. A regra de \texttt{+} vale também para
\texttt{str}, porque os exemplos da especificação concatenam strings, e
\texttt{**} segue a mesma regra de \texttt{-} e \texttt{*}.

\begin{gather*}
\frac{x : \tau \in \Gamma}{\Gamma \vdash x : \tau}
\qquad
\frac{\Gamma \vdash e_1 : \tau \quad \Gamma \vdash e_2 : \tau \quad \tau \in \{\ty{int}, \ty{float}\}}
     {\Gamma \vdash e_1 \oplus e_2 : \tau}\ {\scriptstyle \oplus \in \{+,-,*,**\}}
\\[8pt]
\frac{\Gamma \vdash e_1 : \tau \quad \Gamma \vdash e_2 : \tau \quad \tau \in \{\ty{int}, \ty{float}\}}
     {\Gamma \vdash e_1 \,/\, e_2 : \ty{float}}
\qquad
\frac{\Gamma \vdash e_1 : \ty{int} \quad \Gamma \vdash e_2 : \ty{int}}
     {\Gamma \vdash e_1 \odot e_2 : \ty{int}}\ {\scriptstyle \odot \in \{//, \%\}}
\\[8pt]
\frac{\Gamma \vdash e_1 : \tau \quad \Gamma \vdash e_2 : \tau \quad \tau \in \{\ty{int}, \ty{float}, \ty{str}\}}
     {\Gamma \vdash e_1 < e_2 : \ty{bool}}
\qquad
\frac{\Gamma \vdash e_1 : \ty{list}[\tau] \quad \Gamma \vdash e_2 : \ty{int}}
     {\Gamma \vdash e_1[e_2] : \tau}
\\[8pt]
\frac{\Gamma \vdash f : (\tau_1, \ldots, \tau_n) \to \tau \quad \Gamma \vdash e_i : \sigma_i \quad \sigma_i \le \tau_i}
     {\Gamma \vdash f(e_1, \ldots, e_n) : \tau}
\qquad
\frac{\Gamma \vdash e : C \quad \mathit{campos}(C)(f) = \tau}
     {\Gamma \vdash e.f : \tau}
\\[8pt]
\frac{\texttt{class}\ D(C)}{D \le C}
\qquad
\frac{}{\tau \le \tau}
\qquad
\frac{\tau_1 \le \tau_2 \quad \tau_2 \le \tau_3}{\tau_1 \le \tau_3}
\\[8pt]
\frac{\Gamma, x_1 : \tau_1, \ldots, x_n : \tau_n \vdash e \Leftarrow \tau}
     {\Gamma \vdash \texttt{lambda}\ x_1, \ldots, x_n\texttt{:}\ e \Leftarrow (\tau_1, \ldots, \tau_n) \to \tau}
\end{gather*}

As regras da seção 2.3 do enunciado ficam assim:

\begin{enumerate}\setlength{\itemsep}{0pt}
\item o nome é procurado nos escopos locais, no global e, por último,
  entre as primitivas; se não aparece, o erro é \texttt{undeclared
  variable};
\item operadores conforme as regras acima; \texttt{and}, \texttt{or} e
  \texttt{not} exigem \texttt{bool}, e \texttt{==} e \texttt{!=}, dois
  valores do mesmo tipo (dois objetos, se um for subclasse do outro);
\item na chamada, o número e o tipo dos argumentos; com aridade errada a
  mensagem mostra os dois tipos de função, como no \texttt{.err} do
  enunciado: \texttt{incompatible types: '(int, int) -> int' and
  '(int, int, int) -> ?0'};
\item a expressão do \texttt{return} contra o tipo de retorno, com
  subtipagem;
\item um bloco retorna em todos os caminhos se algum comando dele
  retorna: um \texttt{return}, um \texttt{if} com \texttt{else} em que
  todos os ramos retornam ou um \texttt{while True} sem \texttt{break};
\item \texttt{break} e \texttt{continue} só com o indicador de laço
  ligado, que é desligado na entrada de cada função;
\item nome repetido no escopo corrente, e campo ou método repetido na
  classe;
\item campo e método procurados no ambiente de classes, que já inclui
  os herdados;
\item subtipagem pela cadeia de pais;
\item override com os mesmos tipos de parâmetro e de retorno, exceto o
  \texttt{\_\_init\_\_}, que pode mudar (\texttt{Quadrado} recebe um
  lado, \texttt{Retangulo} recebe dois).
\end{enumerate}

\texttt{len}, \texttt{append}, \texttt{int}, \texttt{float} e
\texttt{str} são casos especiais, com as regras da seção 2.3 do
enunciado; \texttt{print}, \texttt{input}, \texttt{bool} e
\texttt{range} têm tipo fixo, e \texttt{xs.append(e)} vale como
\texttt{append(xs, e)}. Também são verificados \texttt{self} fora de
método (regra 14), o \texttt{\_\_init\_\_} que não atribui algum campo
da classe (regra 13) e a lista vazia sem anotação (seção 6.8).

A subtipagem~\cite{pierce2002} só existe entre classes e entra onde um
valor é usado no lugar de outro: declaração, atribuição, argumento, retorno e elemento de
lista quando o tipo esperado é conhecido. Listas e funções são
invariantes: \texttt{list[Quadrado]} não serve onde se espera
\texttt{list[Forma]}. Num literal de lista ou numa condicional sem tipo
esperado, duas classes diferentes resultam no ancestral comum mais
próximo.

\subsection{Inferência de tipos para MiniPython}

Onde falta anotação o verificador cria uma incógnita (\texttt{TMeta}),
resolvida por unificação~\cite{damas1982}: parâmetro de lambda, lista
vazia, parâmetro ou retorno sem anotação e as variáveis de tipo de uma
função polimórfica a cada chamada. A substituição fica num
\texttt{IntMap}, e ligar \texttt{?0} a um tipo que contém \texttt{?0}
é recusado (occurs check, regra 21). Quando a unificação falha, o erro
mostra o par mais interno que não casou: \texttt{list[int]} contra
\texttt{list[str]} dá \texttt{'int' and 'str'}. Na aridade errada não
há par interno, e saem os dois tipos de função inteiros.

Uma função com variáveis de tipo nas anotações tem um esquema
$\forall a\,b.\,\tau$. Dentro do corpo, \texttt{a} e \texttt{b} são
rígidas (\texttt{TGen}): \texttt{return 1} numa função
\texttt{-> a} é erro. Cada chamada troca as variáveis por incógnitas
novas, e por isso \texttt{identidade(42)} e
\texttt{identidade("olá")} convivem. Uma função de nível superior sem
anotação é verificada com incógnitas, e as que sobram na assinatura, e
não aparecem em outro nome global, viram variáveis de tipo (regra 19);
funções locais não são generalizadas.

A verificação é bidirecional~\cite{dunfield2021}: \texttt{check}
recebe o tipo esperado e o passa para dentro de lambda, literal de lista
e condicional; o resto é inferido e comparado com subtipagem. É assim
que os parâmetros de uma lambda ganham tipo, já que a sintaxe não deixa
anotá-los. Numa chamada, os argumentos que não são lambda são
verificados primeiro. Em \texttt{fold(lambda acc, x: acc + x, 0, nums)},
o \texttt{0} e o \texttt{nums} resolvem \texttt{b = int} e
\texttt{a = int} antes de a lambda ser tipada, e \texttt{acc + x} é
\texttt{int + int}; na ordem do fonte, \texttt{+} receberia duas
incógnitas e não teria como escolher entre \texttt{int}, \texttt{float}
e \texttt{str}.

Uma declaração sem anotação, \texttt{x = e}, recebe o tipo de
\texttt{e}. Se ele ainda tem incógnita, como em \texttt{z = []} ou
\texttt{f = lambda x: x}, o erro é \texttt{cannot infer the type of
'z'} (regra 18).

\subsection{Semântica operacional para MiniPython}

O interpretador segue uma semântica de passos grandes:
\texttt{eval} leva uma expressão e o ambiente a um valor, e
\texttt{exec} executa um comando e diz como ele terminou
(\texttt{Normal}, \texttt{Break}, \texttt{Continue} ou
\texttt{Return v}). Um bloco executa os comandos em ordem até um deles
terminar de outro jeito que não \texttt{Normal}; o laço consome
\texttt{Break} e \texttt{Continue}, e a chamada consome
\texttt{Return}.

\begin{itemize}\setlength{\itemsep}{0pt}
\item \texttt{/} entre inteiros dá \texttt{float}; \texttt{//} e
  \texttt{\%} arredondam para baixo, como no Python (\texttt{div} e
  \texttt{mod} do Haskell); \texttt{and} e \texttt{or} param no
  primeiro operando que decide o resultado.
\item Na chamada, a função e depois os argumentos são avaliados da
  esquerda para a direita; os parâmetros vão para um escopo novo,
  colocado sobre os escopos que a função enxerga: só o global, para
  uma função de nível superior, e ela mesma e o global, para uma
  função local. O resultado é o valor do \texttt{Return}, ou
  \texttt{None}.
\item A lambda guarda a lista de escopos de onde foi criada (a
  closure); na chamada, os parâmetros entram num escopo sobre essa
  lista.
\item \texttt{C(args)} cria um objeto com todos os campos, próprios e
  herdados, valendo \texttt{None}, e chama o \texttt{\_\_init\_\_} da
  classe ou do ancestral mais próximo, com o objeto como
  \texttt{self}.
\item \texttt{obj.m} procura o método a partir da classe concreta do
  objeto, o que dá o despacho dinâmico, e devolve o método já ligado ao
  objeto; \texttt{C.m} procura a partir de \texttt{C}, e é assim que
  \texttt{Retangulo.\_\_init\_\_(self, lado, lado)} funciona.
\item O \texttt{for} percorre a lista pelo índice, relendo-a a cada
  volta, e cada volta abre um escopo com a variável do laço.
\item \texttt{x = e} atualiza \texttt{x} se ele está visível e, se não,
  declara \texttt{x} no escopo mais interno. Dentro de função o global
  não conta para a atribuição, como no Python sem \texttt{global}:
  \texttt{x = e} cria uma local.
\item São erros de execução o índice fora de $[0, \mathit{len}-1]$
  (regra 5; índice negativo não vale), a divisão por zero, a conversão
  de uma string que não é número e o fim da entrada em
  \texttt{input()}.
\end{itemize}

\section{Arquitetura do Compilador}

\subsection{Posição dos nós}

No TP1 a posição de um nó era a do início da construção. Os
\texttt{.err} de \texttt{invalid/semantic} fornecidos com o enunciado
apontam outro lugar: em \texttt{z: float = x + y}, o erro do
\texttt{e01} está na coluna 18, a do \texttt{+}; nos de indexação, no
\texttt{[}; nos de chamada, no \texttt{(}. Por isso operador binário,
chamada, indexação e acesso a campo passaram a guardar a posição do
operador (o \texttt{.} no acesso a campo), nos três parsers; no PEG ela
é capturada com \texttt{@} logo antes do operador. O início da
expressão, que ainda é a posição de um comando e do erro de alvo de
atribuição, sai da função \texttt{exprStart}. Os testes do TP1
continuaram passando, inclusive a comparação das ASTs com posições.

\subsection{Análise semântica}

O verificador (\texttt{Semantic/TypeChecker.hs}) roda numa mônada de
estado com a substituição, o contador de incógnitas, os erros
acumulados, os escopos locais, o global, o ambiente de classes, a
função corrente (tipo de retorno e classe de \texttt{self}) e o
indicador de laço. São duas passadas. Na primeira, as classes e as
funções de nível superior entram no escopo global antes de qualquer
corpo, o que permite recursão e referência mútua; o ambiente de classes
é montado com o pai antes do filho, a herança circular é cortada e
reportada, e o override é conferido aqui. Na segunda, cada item é
verificado na ordem do fonte, de modo que uma função enxerga as
variáveis globais declaradas antes dela.

O verificador não para no primeiro erro. Cada erro é guardado, e a
expressão errada recebe um tipo razoável para a verificação seguir:
um \texttt{\%} com erro dá \texttt{int}, um \texttt{==} dá
\texttt{bool}, um \texttt{+} dá o tipo do operando da esquerda e uma
chamada com aridade errada dá o retorno declarado. Variável não
declarada e indexação inválida dão uma incógnita nova, que unifica com
qualquer coisa e não gera erros em cascata. No \texttt{e01} saem dois
erros, o do \texttt{+} e o da declaração, porque o \texttt{+} com erro
dá \texttt{int}, que não serve para \texttt{float}; no \texttt{e09},
\texttt{[1, 2.0, 3]}, sai um só.

As mensagens estão em inglês, como nos \texttt{.err}, com os tipos
escritos como nas anotações. As incógnitas são renumeradas a partir de
\texttt{?0} em cada mensagem, para que o texto não dependa de quantas
foram criadas antes. Depois dos erros vem uma linha em branco, porque
os \texttt{.err} do enunciado terminam assim. As verificações de
\texttt{return}, \texttt{break} e \texttt{continue} fora de lugar, que
no TP1 ficavam só em \texttt{Validate.hs}, são repetidas no verificador,
para que em \texttt{-{}-check} apareçam junto com os outros erros; nos
modos \texttt{-{}-parse-*} continua valendo o \texttt{Validate}.

\subsection{Interpretador}

A representação dos valores (\texttt{Interpreter/Value.hs}) segue a
sugestão da seção 3.1 do enunciado:

\begin{center}
\begin{tabular}{ll}
\hline
MiniPython & Valor \\
\hline
\texttt{int}, \texttt{float}, \texttt{bool}, \texttt{str} & \texttt{Integer}, \texttt{Double}, \texttt{Bool}, \texttt{String} \\
\texttt{None} & \texttt{VNone} \\
\texttt{list[T]} & \texttt{IORef (Seq Value)} \\
objeto & classe concreta, identificador e \texttt{Map Name (IORef Value)} \\
função \texttt{def} & \texttt{VFun FunDef [Scope]} \\
lambda & \texttt{VLambda [Name] Expr [Scope]} \\
\hline
\end{tabular}
\end{center}

Há ainda \texttt{VPartial}, uma função com os primeiros argumentos já
dados, que é o que \texttt{obj.m} e \texttt{xs.append} devolvem;
\texttt{VClass}, o nome de classe, que chamado cria um objeto; e
\texttt{VBuiltin}, as primitivas. O ambiente é uma lista de escopos,
\texttt{IORef (Map Name (IORef Value))}, o mais interno primeiro e o
global por último. Cada variável tem a sua \texttt{IORef}, então a
atribuição muda o valor no escopo em que a variável foi declarada, e a
closure vê as mudanças feitas depois de criada, como no Python.

Erro de execução é uma exceção com posição, capturada em
\texttt{runProgram}, que devolve \texttt{Either Error ()}; o executável
esvazia a saída padrão antes de imprimir a mensagem. A entrada e a
saída do programa passam por um registro com as ações de ler e de
escrever uma linha: o executável usa \texttt{stdin} e \texttt{stdout},
e o \texttt{cabal test} usa listas de linhas, sem chamar o executável.
\texttt{str} de um \texttt{float} segue o \texttt{repr} do Python:
menor quantidade de dígitos que identifica o número
(\texttt{floatToDigits}) e expoente fora de $[10^{-4}, 10^{16})$, então
\texttt{0.1 + 0.2} dá \texttt{0.30000000000000004} e \texttt{1e16} dá
\texttt{1e+16}.

\section{Resultados e Discussão}

\subsection{Instruções de Uso}

Com GHC 9.6 e cabal: \texttt{cabal update} e \texttt{cabal build}. A
verificação é \texttt{cabal run mpyc -{}- -{}-check exemplos/fib.mpy},
que não imprime nada se o programa está bem tipado, e a execução é
\texttt{cabal run mpyc -{}- -{}-run exemplos/mdc.mpy < exemplos/mdc.in},
que só roda o programa se a verificação passar. Os testes rodam com
\texttt{cabal test} ou \texttt{bash testes/rodar.sh}, e
\texttt{bash testes/gerar-saidas.sh} regenera
\texttt{testes/saidas-tp2.txt}, com código, entrada e saída esperada dos
casos da seção 5 do enunciado.

\subsection{Testes Realizados}

São 106 programas em \texttt{exemplos/}. Os 46 válidos têm que passar
no verificador e ter a saída igual ao \texttt{.out} ao lado, com a
entrada do \texttt{.in} quando o programa lê dados (13 deles). Os 30 de
\texttt{invalid/semantic} (os 15 do enunciado e \texttt{e16} a
\texttt{e30}) têm que dar em \texttt{-{}-check} exatamente o
\texttt{.err}. Os 3 de \texttt{invalid/runtime} passam no verificador e
terminam com o erro de execução do \texttt{.err}, e os 27 de erro léxico
ou sintático vêm do TP1. Os 106 passam nas duas suítes. Os casos pedidos
na seção 5 do enunciado são:

\begin{center}
\footnotesize
\begin{tabular}{ll|ll}
\hline
Caso válido & Programa & Caso de erro & Programa \\
\hline
MDC & \texttt{secao\_9\_01\_mdc} & variável não declarada & \texttt{e05} \\
Fibonacci & \texttt{secao\_9\_02\_fibonacci} & \texttt{int + str} & \texttt{e16} \\
busca linear & \texttt{secao\_9\_03\_busca\_linear} & aridade errada & \texttt{e04} \\
seleção & \texttt{secao\_9\_04\_selecao} & falta de \texttt{return} & \texttt{e03} \\
\texttt{map} e \texttt{filter} & \texttt{secao\_9\_08\_alta\_ordem} & \texttt{break}/\texttt{continue} fora de laço & \texttt{e17} \\
formas & \texttt{secao\_9\_11\_formas} & campo inexistente & \texttt{e18} \\
pilha & \texttt{secao\_9\_12\_pilha} & & \\
lambdas & \texttt{secao\_9\_15\_lambdas} & & \\
\hline
\end{tabular}
\end{center}

Os outros programas de erro cobrem as demais regras (método
inexistente, override, um \texttt{Animal} onde se espera
\texttt{Cachorro}, redeclaração, escopo de bloco, campo não
inicializado, \texttt{self} fora de método, entre outras), e um deles
tem cinco erros, para mostrar o acúmulo:

\begin{lstlisting}[style=gram]
e26_varios_erros.mpy:2:5: error: incompatible types: 'float' and 'int'
e26_varios_erros.mpy:5:17: error: operator '/' requires operands of the same type (int or float), got 'float' and 'int'
e26_varios_erros.mpy:8:26: error: incompatible types: 'float' and 'str'
e26_varios_erros.mpy:9:21: error: operator '+' requires operands of the same type (int or float), got 'str' and 'float'
e26_varios_erros.mpy:10:15: error: undeclared variable: 'total'
\end{lstlisting}

As saídas dos válidos foram conferidas com o Python 3: os mesmos
programas, sem as anotações de tipo (\texttt{(int) -> int} não é
sintaxe Python) e com \texttt{append} e \texttt{range} definidos,
deram as 46 saídas iguais. Os \texttt{.err} do enunciado passam byte a
byte, menos o do \texttt{e03}, que veio vazio; como o TP2 pede a regra
5, ele agora tem \texttt{missing return in function 'f'}.

\subsection{Dificuldades encontradas}

\paragraph{Posição dos erros.} A convenção do TP1, posição no início
da construção, foi escolhida pensando nos erros de tipo, mas os
\texttt{.err} do enunciado apontam o operador. Foi preciso voltar
atrás nos três parsers, guardar a posição do operador e criar
\texttt{exprStart} para os lugares que usavam o início.

\paragraph{Unificação e subtipagem.} Unificar o tipo da função com o
tipo dos argumentos, como sugerem os \texttt{.err}, recusaria
\texttt{imprimir(r)} com \texttt{r} do tipo \texttt{Retangulo} e o
parâmetro do tipo \texttt{Forma}. A subtipagem ficou só nos pontos em
que um valor é usado no lugar de outro, comparando classes pela cadeia
de pais, e a unificação no resto.

\paragraph{Lambdas sem anotação.} Inferida sozinha, uma lambda como
\texttt{lambda acc, x: acc + x} tem dois parâmetros de tipo
desconhecido, e \texttt{+} não sabe qual regra aplicar. A solução foi a
verificação bidirecional com os argumentos que não são lambda
verificados antes.

\paragraph{Erros em cascata.} Na primeira versão,
\texttt{self.valor = self.valor + 1} numa função comum gerava cinco
mensagens. O \texttt{self} fora de método passou a ser um parâmetro de
tipo desconhecido, com um erro só, e o acesso a campo de um valor de
tipo desconhecido, que já vem de um erro, deixou de gerar mensagem.

\paragraph{Variável global dentro de função.} A primeira versão deixava
a função reatribuir uma global. A comparação com o Python mostrou a
diferença: lá, \texttt{contador = 5} dentro da função cria uma local, e
\texttt{g = g + 1} é erro. Como o enunciado manda adotar o Python
quando a especificação é omissa, \texttt{x = e} passou a criar uma
local e \texttt{x += e} numa global passou a ser erro de compilação.

\paragraph{Formato dos \texttt{.err}.} Com as mensagens já iguais, o
\texttt{cabal test} ainda falhava nos 14 \texttt{.err} do enunciado: os
arquivos terminam com uma linha em branco a mais, que o
\texttt{-{}-check} passou a imprimir.

\subsection{Limitações}

A lambda enxerga todos os escopos de fora, enquanto a especificação
limita a captura ao escopo imediatamente envolvente e ao global. Listas
e funções são invariantes, sem a contravariância dos parâmetros. A
regra 13 só confere se cada campo é atribuído em algum ponto do
\texttt{\_\_init\_\_}, e não em todos os caminhos nem antes do uso.
Duas funções sem anotação que se chamam uma à outra ficam monomórficas,
porque a generalização é feita depois de verificar cada uma.
\texttt{int}, \texttt{float} e \texttt{str} só podem ser chamadas, e não
usadas como valor, porque aceitam mais de um tipo de argumento.

\section{Uso de IA generativa}

Como no TP1, foi usado o assistente Claude (Anthropic), agora pelo
Claude Code. Ele escreveu o código do verificador, do interpretador e
dos testes e o rascunho deste relatório; os prompts, as respostas e as
modificações (item 4 da seção 6 do enunciado) estão no
Apêndice~\ref{sec:ia}, fora do limite de oito páginas.

\section{Conclusão}

O verificador confere as dez regras pedidas e reproduz os
\texttt{.err} do enunciado, e o interpretador dá as mesmas saídas do
Python. O mais trabalhoso foi juntar unificação, necessária para
lambdas e funções polimórficas, com subtipagem entre classes.

\begin{thebibliography}{9}
\setlength{\itemsep}{0pt}
\bibitem{pierce2002}
Pierce, B. C. (2002). \emph{Types and Programming Languages}. MIT Press.

\bibitem{damas1982}
Damas, L.; Milner, R. (1982). Principal type-schemes for functional programs. In \emph{POPL '82}. ACM.

\bibitem{dunfield2021}
Dunfield, J.; Krishnaswami, N. (2021). Bidirectional Typing. \emph{ACM Computing Surveys}, 54(5).
\end{thebibliography}

\clearpage
\appendix
% Apendice do relatorio do TP2: uso de IA generativa (item 4 da
% secao 6 do enunciado). Fica fora do limite de oito paginas.

\section{Uso de IA generativa no TP2}
\label{sec:ia}

A ferramenta foi o Claude Code (Anthropic), com o modelo Claude Opus
5.5, numa única sessão em 25/11/2026, aberta no diretório do
repositório. Nela a ferramenta lia arquivos, rodava comandos
(\texttt{cabal}, os testes, o Python) e editava o código. Os prompts
estão abaixo como foram digitados, sem correção. Os três últimos
respondem a perguntas que a ferramenta fez ao fim da resposta anterior.

\subsection*{Prompt 1}

\begin{quote}
vamos continuar esse trabalho analise o repositorio veja oq da pra
melhorar e fazer as modificações necessarias para o TP2\\
\url{https://github.com/Brendagnasc/bcc328-minipython}\\
\emph{[texto completo do enunciado do TP2, colado]}\\
Repositorio do professor\\
\url{https://github.com/rodrigogribeiro/bcc328-2026-2}
\end{quote}

\textbf{Resposta.} A ferramenta leu o repositório, a especificação e o
repositório do professor, que não tinha material do TP2 além das notas
de aula. Concluiu que os \texttt{.err} de \texttt{invalid/semantic},
fornecidos com o enunciado do TP1, vinham de um verificador de
referência e serviam de modelo para as mensagens e as posições. Em
seguida:

\begin{enumerate}\setlength{\itemsep}{0pt}
\item mudou a posição guardada nos nós de operador, chamada, indexação
  e acesso a campo nos três parsers, criou \texttt{exprStart} e
  confirmou que os 88 testes do TP1 continuavam passando;
\item escreveu \texttt{Semantic/Types.hs} e
  \texttt{Semantic/TypeChecker.hs}, com as decisões descritas na
  Seção 2, e \texttt{Interpreter/Value.hs} e
  \texttt{Interpreter/Interp.hs};
\item acrescentou \texttt{-{}-check} e \texttt{-{}-run} em
  \texttt{CLI.hs} e \texttt{app/Main.hs};
\item comparou a saída de \texttt{-{}-check} com os 15 \texttt{.err}
  do enunciado: 14 iguais, e o \texttt{e03} diferente por estar vazio;
\item escreveu, fora do repositório, um tradutor de MiniPython para
  Python que tira as anotações de tipo, e comparou as saídas dos 46
  programas válidos com as do Python 3;
\item criou os programas \texttt{e16} a \texttt{e30}, os de
  \texttt{invalid/runtime} e os arquivos \texttt{.in} e \texttt{.out},
  atualizou as duas suítes de testes e o \emph{script} que gera as
  saídas esperadas, e escreveu o README e o CHANGELOG do TP2.
\end{enumerate}

\textbf{Modificações.} Durante a própria sessão, estes pontos da
resposta foram corrigidos depois de compilar ou testar:

\begin{itemize}\setlength{\itemsep}{0pt}
\item duas partes do interpretador foram reescritas antes de compilar:
  a atribuição aumentada num elemento de lista, cuja primeira versão
  não gravava o valor novo de volta na lista, e a função de
  comparação, que usava um tipo que precisaria de uma extensão do GHC
  não habilitada;
\item a primeira versão deixava uma função reatribuir uma variável
  global; a comparação com o Python mostrou que lá isso cria uma
  variável local, e a regra foi trocada pela do Python;
\item \texttt{z = []} sem anotação era aceito e passou a ser erro, como
  pede a seção 6.8 da especificação, e a mensagem
  \texttt{'int', 'float', or 'bool'} perdeu a vírgula antes do
  \texttt{or};
\item um \texttt{self} usado fora de método gerava cinco mensagens
  numa linha, e passou a gerar uma;
\item com as mensagens já iguais, o \texttt{cabal test} falhou nos 14
  \texttt{.err} do enunciado, porque eles terminam com uma linha em
  branco; o \texttt{-{}-check} passou a imprimi-la;
\item o \texttt{e03\_funcao\_sem\_return.err}, vazio, recebeu a
  mensagem do verificador, porque o TP2 exige a regra 5.
\end{itemize}

\subsection*{Prompt 2}

\begin{quote}
faca os commits mas sem coautor do claude
\end{quote}

\textbf{Resposta.} Foram feitos seis commits, em nome da autora e sem a
linha de coautoria que a ferramenta coloca por padrão: posição do
operador nos parsers, verificador de tipos, interpretador, opções da
linha de comando, testes e documentação. O \texttt{minipython.cabal},
que tinha trechos de quatro desses commits, foi dividido entre eles.
Cada commit foi compilado com \texttt{-Werror} e testado num diretório
de trabalho separado: os quatro primeiros passaram nos 88 testes do
TP1 e os dois últimos nos 106. Não houve modificação.

\subsection*{Prompt 3}

\begin{quote}
Sim, dar push
\end{quote}

\textbf{Resposta.} Os seis commits foram enviados para o GitHub. Não
houve modificação.

\subsection*{Prompt 4}

\begin{quote}
Sim, escrever relatório
\end{quote}

\textbf{Resposta.} A ferramenta escreveu este relatório e este
apêndice, em arquivos separados do relatório do TP1, a partir do
código, dos testes e do histórico da sessão, e compilou o PDF
conferindo o limite de oito páginas.

\subsection*{Avaliação}

O código veio da ferramenta, mas a maior parte dos problemas foi
encontrada pelos testes: os \texttt{.err} do enunciado, que mostraram a
questão da posição e da linha em branco, e a comparação com o Python,
que mostrou a regra das variáveis globais. Onde a especificação é
omissa a ferramenta tomou decisões (escopo por bloco, \texttt{+} entre
strings, \texttt{==} para qualquer tipo, a regra das globais), e todas
estão descritas na Seção 2 e no README, para que possam ser avaliadas.


\end{document}
